\section{Algoritmet e diferencimit dhe kuadraturës}
\subsection{Derivati qendror dhe studimi i hapit}
\begin{lstlisting}[language=Python,caption={Gabimi i derivatit qendror për një varg hapash.}]
import numpy as np

def derivat_qendror(f, x, h):
    return (f(x + h) - f(x - h)) / (2.0*h)

f = np.sin
x0 = 1.0
h = np.logspace(-16, -1, 200)
D = derivat_qendror(f, x0, h)
gabimi = np.abs(D - np.cos(x0))
h_opt = h[np.argmin(gabimi)]
print("h optimal empirik =", h_opt)
\end{lstlisting}
Grafiku $E(h)$ duhet të paraqitet në bosht log--log. Në zonën e këputjes pritet pjerrësi afërsisht 2; në zonën e rrumbullakimit gabimi rritet kur $h$ zvogëlohet.

\subsection{Rregulli i përbërë i Simpsonit}
\begin{lstlisting}[language=Python,caption={Integrimi Simpson me kontrollin e numrit çift të nënintervaleve.}]
def simpson(f, a, b, n):
    if n <= 0 or n % 2 != 0:
        raise ValueError("n duhet te jete numer pozitiv çift.")
    x = np.linspace(a, b, n + 1)
    y = f(x)
    h = (b - a) / n
    return h/3 * (y[0] + y[-1]
                  + 4*np.sum(y[1:-1:2])
                  + 2*np.sum(y[2:-2:2]))

for n in [4, 8, 16, 32, 64]:
    I = simpson(np.sin, 0.0, np.pi, n)
    print(n, I, abs(I - 2.0))
\end{lstlisting}

\subsection{Vlerësimi Runge i gabimit}
Nëse metoda ka rend $p$, dy llogaritje me $h$ dhe $h/2$ japin
\begin{equation}
 E_{h/2}\approx\frac{|A(h/2)-A(h)|}{2^p-1}.
\end{equation}
Ky vlerësim është i vlefshëm vetëm në regjimin asimptotik. Prandaj kontrollohet edhe raporti i gabimeve për tri hapa radhazi.

\begin{lstlisting}[language=Python]
def studim_konvergjence(llogarit, ns, p):
    v = np.array([llogarit(n) for n in ns])
    e_runge = np.abs(v[1:] - v[:-1]) / (2**p - 1)
    return v, e_runge
\end{lstlisting}

\subsection{Kuadratura e Gausit}
\begin{lstlisting}[language=Python,caption={Gauss--Legendre në një interval të përgjithshëm.}]
from numpy.polynomial.legendre import leggauss

def gauss_legendre(f, a, b, n):
    xi, w = leggauss(n)
    x = 0.5*(b-a)*xi + 0.5*(a+b)
    return 0.5*(b-a) * np.sum(w * f(x))
\end{lstlisting}
Në notebook krahasohen numri i vlerësimeve të funksionit, gabimi dhe koha e ekzekutimit. Kjo e lidh kuadraturën klasike me motivimin për Monte Carlo në dimension të lartë.

Formulat dhe emërtimet shqip të diferencimit e integrimit numerik mbështeten te \cite{hanelli,hoxha,datja}; interpretimi fizik dhe krahasimi me Monte Carlo plotësohen nga \cite{newman,press}.
